\documentclass[./examen.tex]{subfiles}

\begin{document}
\begin{center}
\makebox[\textwidth]{Nombre y Apellido:\enspace\hrulefill}
\end{center}

\begin{questions}
\question ¿A qué tensión debemos conectar un condensador de $C=104\mu F$ para que adquiera una carga de $Q=2x10^{-6}C$?
\question Calcular la carga de un condensador de $C=200\mu F$ si se conecta a $V=100V$ y su capacidad si se conectan 4 en serie.
       \question Calcular las permitividades absolutas teniendo las relativas $\epsilon_{vidrio}=4.7,\quad \epsilon_{mica}=5.6\quad y\quad  \epsilon_{glicerina}=45.$
       
       
    \question Una bobina de $N=200$ espiras es recorrida por una corriente $I=20A$ que da lugar a un flujo magnético de $\Phi=0,002Wb$. Calcular el coeficiente de autoinducción $L$ de la bobina.

     \question Se desea saber capacidad en faradios de un capacitor que se carga con $Q=1,8C$ y posee placas paralelas de $S=1m^2$ si el dieléctrico es el glicerina y la distancia entre placas de 0.001mm.
       
\end{questions}

 \end{document}